% Generated draft: requires mathematical and source review.
\documentclass[11pt]{article}
\usepackage[margin=25mm]{geometry}
\usepackage{fontspec}
\setmainfont{Latin Modern Roman}
\usepackage{amsmath,amssymb,mathtools,unicode-math}
\setmathfont{Latin Modern Math}
\usepackage{enumitem,needspace,xurl,hyperref}
\setlength{\parindent}{0pt}
\setlength{\parskip}{5pt}
\setlength{\emergencystretch}{3em}
\newcommand{\N}{\mathbb{N}}
\newcommand{\Z}{\mathbb{Z}}
\newcommand{\Q}{\mathbb{Q}}
\newcommand{\R}{\mathbb{R}}
\newcommand{\C}{\mathbb{C}}
\newcommand{\F}{\mathbb{F}}
\newcommand{\A}{\mathbb{A}}
\newcommand{\PP}{\mathbb{P}}
\newcommand{\E}{\mathbb{E}}
\newcommand{\eps}{\varepsilon}
\newcommand{\dd}{\,\mathrm{d}}
\begin{document}
\section*{UnsolvedMath ID: 11100036}
\textbf{Source problem code:} catalogue entry
\subsection{Definitions and mathematical statement}
For a commutative ring $R$, let $D(R)$ be its unbounded derived category and $\operatorname{Perf}(R)$ its perfect complexes, represented by bounded complexes of finitely generated projective modules. A thick subcategory is triangulated and closed under retracts; a localizing subcategory is triangulated and closed under arbitrary coproducts. A tensor ideal is additionally closed under tensoring by objects of the ambient category.

A Thomason subset of $\operatorname{Spec}R$ is a union $\bigcup_i Z_i$ of closed subsets whose complements are quasi-compact. The source recalls Thomason's classification of thick subcategories of finite objects by these subsets. In the affine commutative case the tensor-ideal condition is automatic; for a general scheme the standard statement concerns thick tensor ideals of perfect complexes, with scheme hypotheses, and does not classify all thick subcategories unqualified.

Axiomatic stable homotopy item 4 asks for the corresponding study/classification of localizing subcategories and for extensions of Thomason's methods to more general non-Noetherian stable homotopy categories. It does not assert that every such category has a classification by arbitrary prime-spectrum subsets. For schemes one must specify a category such as $D_{\mathrm{qc}}(X)$, rather than silently identifying it with the derived category of all sheaves, and distinguish localizing subcategories from tensor-ideal ones.

The registry reports the commutative Noetherian ring case classified by arbitrary subsets of $\operatorname{Spec}R$ (Neeman), and attributes thick perfect-complex classification to Thomason, with Balmer-spectrum context. It mentions further stratification work for schemes and unresolved non-Noetherian scope, but its final sentence is truncated. No precise general non-Noetherian category, tensor structure, coefficient-ring action, finiteness axioms or target spectrum is supplied, so those must be part of a specified extension rather than guessed.
\subsection{Short English statement}
Extend the known perfect-complex support classification to localizing subcategories and broader non-Noetherian stable categories. Scheme, tensor-ideal and derived-category conventions materially affect the question.
\subsection{Sources}
\begin{itemize}
\item Local registry: \url{data/unsolvedmath/problems.json}, permanent UnsolvedMath ID $11100036$, exact source code \texttt{catalogue entry}, attributed to Mark Hovey.
\item \emph{Hovey's Algebraic Topology Problem List}, Axiomatic stable homotopy item 4: \url{https://www-users.cse.umn.edu/~tlawson/hovey/}. Supplied access date: 2026-07-29; extraction: html-ordered-list; 65 top-level items in nine category pages.
\item The source invokes one of Bob Thomason's last papers without a full citation. The registry additionally names Neeman, Balmer, Antieau--Heller, Barthel--Heard--Sanders, and Dell'Ambrogio--Stevenson and others for derived-category/stratification work, but supplies no precise papers or hypotheses; its last progress sentence is truncated.
\end{itemize}
\end{document}
